\section{Ingeniería del harness: codiseño de herramientas, interfaz y flujo de control}

\subsection{Por qué el harness es un ciudadano de primera clase}
Sutton dice a los 00:26 una frase que se cita con frecuencia en esta clase: \textit{``The design of the agent harness and the tools, you have to co-design it with whatever the model can do.''}
Esta frase significa que no debemos tratar al modelo como una caja negra a la que se le conectan herramientas por fuera, sino tratar la interfaz de las herramientas, el formato del prompt y la estructura del flujo de control como objetos de diseño del mismo nivel que el modelo.

\begin{importantbox}{Tres principios de ingeniería del Harness Engineering}
\begin{enumerate}
    \item La semántica de las acciones debe acercarse a la distribución que el modelo conoce, para reducir la ambigüedad en las llamadas;
    \item La granularidad de las herramientas debe ser lo bastante compacta para reducir los casos de ``llamar a cinco herramientas para hacer una sola cosa'';
    \item El formato de salida debe ser comprimible y rastreable, para evitar que el ruido sature el contexto.
\end{enumerate}
\end{importantbox}

\subsection{La perspectiva de Agent Computer Interface (ACI)}
En la clase se cita el concepto de ``ACI'' (Agent Computer Interface), con la intención de establecer una analogía con HCI: el diseño de interfaces persona-computadora tiene paradigmas establecidos, y la interfaz entre el agente y su entorno de ejecución necesita, de igual manera, un método sistemático. Por ejemplo, acciones como ``ir a la definición'' o ``buscar referencias por símbolo'' suelen ajustarse mejor a la ruta de pensamiento del modelo que un grep genérico, y se acercan más al flujo de trabajo de un IDE.

\begin{knowledgebox}{El principio de la ``ruta familiar'' en el diseño de herramientas}
El presentador retoma una idea parecida al ``Little Red Riding Hood principle'': conviene llevar al modelo, en la medida de lo posible, hacia los numerosos patrones de interacción que ya vio durante el entrenamiento posterior. Cuanto más se acerquen el nombre de la herramienta, sus parámetros y el texto que devuelve a los contextos de desarrollo habituales, más fácil será lograr llamadas exitosas de manera estable.
\end{knowledgebox}

\begin{figure}[H]
\centering
\includegraphics[width=0.88\textwidth]{frame-03-harness.jpg}
\caption{Imagen clave de la clase que enfatiza que ``el harness y las herramientas deben diseñarse en conjunto'' \protect\footnotemark}
\end{figure}
\footnotetext{Intervalo de tiempo en el video: 00:26:24--00:26:55.}

\subsection{El espectro del diseño del flujo de control: de la libertad a la restricción}
La clase ubica el flujo de control de los agentes de programación en un espectro continuo: en un extremo está ``dar shell más IDE, casi todo dinámico''; en el medio están los ciclos ReAct como el de SWE-agent; más hacia la derecha están los ensambles por etapas (primero recuperación de información, luego generación del parche); y en el extremo derecho están las máquinas de estados (state machine) explícitas, que restringen las acciones ejecutables en cada paso.

\begin{table}[H]
\centering
\begin{tabular}{p{3.2cm}p{4.4cm}p{3.8cm}p{2.6cm}}
\toprule
\textbf{Paradigma de flujo de control} & \textbf{Ventajas} & \textbf{Riesgos} & \textbf{Tareas idóneas} \\
\midrule
ReAct totalmente dinámico & Exploración flexible, se adapta a situaciones desconocidas & Deriva de la trayectoria, costo inestable & Localización y depuración abiertas \\
Ciclo de dos etapas ensambladas & Desacopla la recolección de información de la generación del parche & Un error en el límite entre etapas provoca fallas en cascada & Tareas de reparación de complejidad media \\
Ciclo restringido por máquina de estados & Alta controlabilidad, facilita la auditoría y la reproducción & Espacio de búsqueda limitado, puede pasar por alto soluciones & Escenarios de alto riesgo, procesos de cumplimiento \\
\bottomrule
\end{tabular}
\caption{El diseño del flujo de control es un equilibrio entre el riesgo de la tarea y la capacidad de exploración}
\end{table}

\begin{warningbox}{``Mientras más control, mejor'' es un malentendido}
Una restricción demasiado fuerte le quita al Agent la capacidad de recuperarse en rutas anómalas; una restricción demasiado débil produce alucinaciones en trayectorias largas. Lo correcto es configurar el flujo de control por capas, según la complejidad y el nivel de riesgo de la tarea, en lugar de buscar un solo paradigma.
\end{warningbox}

\subsection{Análisis de trayectorias: hay que revisar las muestras fallidas}
Antes de abordar el tema de los agentes de seguridad, Sutton repite un viejo principio del aprendizaje automático: \textit{``Always look at the data.''} Aquí ``data'' no se refiere al conjunto de entrenamiento, sino a las trayectorias de ejecución del Agent. Solo revisando una por una las rutas fallidas se puede saber si el problema fue un desajuste en la recuperación, un mal uso de una herramienta o un verificador que indujo a error.

\begin{importantbox}{Recomendación práctica: iteración del harness guiada por trayectorias}
\begin{itemize}
    \item En cada iteración de versión, tomar una muestra fija de trayectorias fallidas para revisarlas de manera manual;
    \item Traducir los patrones de falla más frecuentes en nuevas herramientas o nuevas restricciones;
    \item Consolidar la experiencia de ``ya lo arreglé una vez'' en la system instruction y en las plantillas de política.
\end{itemize}
\end{importantbox}

\begin{figure}[H]
\centering
\includegraphics[width=0.88\textwidth]{frame-04-control-loop.jpg}
\caption{Explicación centrada en el ciclo de flujo de control y de política, que enfatiza que el ciclo dinámico y la restricción programática no son opuestos \protect\footnotemark}
\end{figure}
\footnotetext{Intervalo de tiempo en el video: 00:38:10--00:38:55.}

\subsection{Resumen de la sección}
El núcleo de Harness Engineering es la ``coordinación del sistema'', no el ``ajuste fino del prompt''. La semántica de las herramientas, la estructura de la retroalimentación, las restricciones del flujo de control y la capacidad del modelo deben optimizarse en conjunto para poder mejorar de manera estable en tareas complejas.
